\bta{2020}

\begin{enumerate}
	%\renewcommand{\labelenumi}{\arabic{enumi}.}
	% A(\Alph) a(\alph) I(\Roman) i(\roman) 1(\arabic)
	%设定全局标号series=example	%引用全局变量resume=example
	%[topsep=-0.3em,parsep=-0.3em,itemsep=-0.3em,partopsep=-0.3em]
	%可使用leftmargin调整列表环境左边的空白长度 [leftmargin=0em]
	\item
	（共 30 分）考虑一维束缚态。
	\begin{enumerate}
		%\renewcommand{\labelenumi}{\arabic{enumi}.}
		% A(\Alph) a(\alph) I(\Roman) i(\roman) 1(\arabic)
		%设定全局标号series=example	%引用全局变量resume=example
		%[topsep=-0.3em,parsep=-0.3em,itemsep=-0.3em,partopsep=-0.3em]
		%可使用leftmargin调整列表环境左边的空白长度 [leftmargin=0em]
		\item
		证明 $\langle\psi(x, t) \mid \psi(x, t)\rangle$ 不随时间变化, 此处波函数 $\psi$ 不必是定态。
		\item 
		证明对于定态, 动量的期望值为零。
		\item 
		证明如果粒子在 $t=0$ 时刻处于定态, 则在以后时刻永远保持定态。
		
		
		
	\end{enumerate}
\banswer{
	\begin{enumerate}
		\item 由薛定谔方程及其伴随式，
		\begin{align*}
			\frac{d}{dt}\langle\psi|\psi\rangle
			&=\langle\dot\psi|\psi\rangle+\langle\psi|\dot\psi\rangle\\
			&=\frac{i}{\hbar}\langle\psi|H|\psi\rangle
			-\frac{i}{\hbar}\langle\psi|H|\psi\rangle=0,
		\end{align*}
		其中用了 $H^\dagger=H$。故任意态的范数均守恒。

		\item 一维实势的束缚能级非简并，定态空间波函数可取为实函数。于是
		\[
			\langle p\rangle=-i\hbar\int_{-\infty}^{\infty}
			\psi(x)\psi'(x)dx
			=-\frac{i\hbar}{2}\left[\psi^2(x)\right]_{-\infty}^{\infty}=0.
		\]
		束缚态在无穷远处消失，故边界项为零。

		\item 若 $H|\phi_n\rangle=E_n|\phi_n\rangle$ 且
		$|\psi(0)\rangle=|\phi_n\rangle$，则
		\[
			\boxed{|\psi(t)\rangle=e^{-iE_nt/\hbar}|\phi_n\rangle}.
		\]
		它只多一个整体相位，仍是同一本征值 $E_n$ 的定态。
	\end{enumerate}
}
	
\item 
（共 30 分）设波函数$\psi(x)=(1 / \sqrt{2 \pi}) e^{i p(x+\beta) \hbar}$, 而 $\hat{x}, \hat{p}$ 分别为 $x$ 方向的坐标
和动量算符，其中 $\beta$ 为实常数。
\begin{enumerate}
	%\renewcommand{\labelenumi}{\arabic{enumi}.}
	% A(\Alph) a(\alph) I(\Roman) i(\roman) 1(\arabic)
	%设定全局标号series=example	%引用全局变量resume=example
	%[topsep=-0.3em,parsep=-0.3em,itemsep=-0.3em,partopsep=-0.3em]
	%可使用leftmargin调整列表环境左边的空白长度 [leftmargin=0em]
	\item
说明 $\psi(x)$ 是否为 $\hat{p}$ 的归一化本征态。
\item 
证明 $\left\langle x^{\prime}\right| e^{i \alpha \hat{p} / \hbar}=\left\langle x^{\prime}+\alpha\right|$, 及 $e^{i \alpha \hat{p} / \hbar}\left|x^{\prime}\right\rangle=\left|x^{\prime}-\alpha\right\rangle$, 其中 $\alpha$ 为实常数。
\item 
化简算符 $e^{i \alpha \hat{p} / \hbar} \hat{x} e^{-i \alpha \hat{p} / \hbar}$ 。
\item 
化简算符 $e^{i \alpha \hat{p} / \hbar} \hat{x}^{2} e^{-i \alpha \hat{p} / \hbar}$ 。

\end{enumerate}

\banswer{
	题中指数应按通常动量本征函数理解为
	$\exp[ip(x+\beta)/\hbar]$；$\beta$ 只产生整体相位。
	\begin{enumerate}
		\item
		\[
			-i\hbar\frac{d}{dx}\psi(x)=p\psi(x),
		\]
		所以它是动量本征值 $p$ 的广义本征态。但平面波在整条实轴上不可平方归一：
		$\int|\psi|^2dx=\infty$；它只能按 Dirac $\delta$ 函数归一化。严格的标准系数应随采用
		$\delta(p-p')$ 还是 $\delta(k-k')$ 归一化而取
		$1/\sqrt{2\pi\hbar}$ 或 $1/\sqrt{2\pi}$。

		\item 记 $T(\alpha)=e^{i\alpha p/\hbar}$。由下一问的共轭关系或直接在坐标表象作用可得
		\[
			\langle x'|T(\alpha)|\psi\rangle=\psi(x'+\alpha)
			=\langle x'+\alpha|\psi\rangle,
		\]
		故
		\[
			\boxed{\langle x'|T(\alpha)=\langle x'+\alpha|}.
		\]
		取伴随或投影到任意坐标态，得到
		\[
			\boxed{T(\alpha)|x'\rangle=|x'-\alpha\rangle}.
		\]

		\item 用 Baker--Campbell--Hausdorff 公式及 $[p,x]=-i\hbar$：
		\[
			\boxed{e^{i\alpha p/\hbar}x e^{-i\alpha p/\hbar}
			=x+\alpha}.
		\]
		更高重对易子为零。

		\item 幺正共轭保持乘积，故
		\[
			\boxed{e^{i\alpha p/\hbar}x^2e^{-i\alpha p/\hbar}
			=(x+\alpha)^2}.
		\]
	\end{enumerate}
}


\item 
（ 共 30 分 ） 一 个 无 自 旋 粒 子 的 波 函 数 为
$\psi=K(x+i y+2 z) e^{-\alpha r} $， 此 处
$r=\sqrt{x^{2}+y^{2}+z^{2}}$, 其中 $K, \alpha$ 为实常数。  （球谐函数: $Y_{0}^{0}=\sqrt{\frac{1}{4 \pi}}$, $ Y_{1}^{0}=\sqrt{\frac{3}{4 \pi}} \cos \theta$, 
$Y_{1}^{\pm 1}=\mp \sqrt{\frac{3}{8 \pi}} \sin \theta e^{\pm i \phi} $）
\begin{enumerate}
	%\renewcommand{\labelenumi}{\arabic{enumi}.}
	% A(\Alph) a(\alph) I(\Roman) i(\roman) 1(\arabic)
	%设定全局标号series=example	%引用全局变量resume=example
	%[topsep=-0.3em,parsep=-0.3em,itemsep=-0.3em,partopsep=-0.3em]
	%可使用leftmargin调整列表环境左边的空白长度 [leftmargin=0em]
	\item
	求粒子的总角动量。
	\item 
	求角动量 $z$ 分量即 $\hat{L}_{z}$ 的期望值, 及测得 $L_{z}=\hbar$ 的概率。
	\item 
	求发现粒子在 $(\theta, \varphi)$ 方向上 $d \Omega$ 立体角内的概率。
	
	
	
\end{enumerate}

\banswer{
	写成球坐标后，角向因子为
	\[
		x+iy+2z=r\left(\sin\theta e^{i\varphi}+2\cos\theta\right)
		=r\left[-\sqrt{\frac{8\pi}{3}}Y_1^1
		+2\sqrt{\frac{4\pi}{3}}Y_1^0\right].
	\]
	\begin{enumerate}
		\item 波函数完全位于 $l=1$ 子空间，故
		\[
			\boxed{L^2=2\hbar^2},\qquad |\vec L|=\sqrt2\,\hbar.
		\]

		\item $m=1$ 与 $m=0$ 分量的权重之比为 $1:2$，所以
		\[
			\boxed{\langle L_z\rangle=\frac13\hbar},\qquad
			\boxed{P(L_z=\hbar)=\frac13}.
		\]

		\item 归一化角分布为
		\[
			w(\theta,\varphi)=\frac1{8\pi}
			\left|\sin\theta e^{i\varphi}+2\cos\theta\right|^2.
		\]
		因此在指定方向的 $d\Omega$ 内发现粒子的概率为
		\[
			\boxed{dP=\frac{\sin^2\theta+4\cos^2\theta
			+4\sin\theta\cos\theta\cos\varphi}{8\pi}\,d\Omega}.
		\]
		径向部分已积分掉；若需要全波函数归一化，则
		$K=\sqrt{\alpha^5/(6\pi)}$。
	\end{enumerate}
}


\newpage
\item 
（共 30 分）
\begin{enumerate}
	%\renewcommand{\labelenumi}{\arabic{enumi}.}
	% A(\Alph) a(\alph) I(\Roman) i(\roman) 1(\arabic)
	%设定全局标号series=example	%引用全局变量resume=example
	%[topsep=-0.3em,parsep=-0.3em,itemsep=-0.3em,partopsep=-0.3em]
	%可使用leftmargin调整列表环境左边的空白长度 [leftmargin=0em]
	\item
	一个电子在 $t=0$ 的时刻处于自旋态 $\chi=\frac{1}{3}\left(\begin{array}{c}1-2 i \\ 2\end{array}\right)$ 。在 $t>0$ 时刻，在外界加一个
	磁场 $\vec{B}=B_{0}\left(\sin \theta \hat{e}_{x}+\cos \theta \hat{e}_{z}\right)$, 此时电子的哈密顿量为 $\hat{H}=-2 \mu_{B} \hat{\vec{S}} \bullet \vec{B}$, 其中 $\hat{\vec{S}}$ 为 自旋算符, $\mu_{B}$ 为玻尔磁子，求此粒子在任意 $t$ 时刻的波函数。
	\item 
	考虑两个自旋为 $\frac{1}{2}$ 的粒子处于磁场中，此时系统的哈密顿量为
	\[ \hat{H}=a \hat{\sigma}_{1 z}+b \hat{\sigma}_{2 z}+c_{0} \hat{\bar{\sigma}}_{1} \bullet \hat{\vec{\sigma}}_{2} \]
	其中 $a, b, c_{0}$ 为常数, $\hat{\vec{\sigma}}$是泡利算符，前两项为粒子处于磁场中的势能，最后
	一项为两粒子自旋-自旋相互作用能。求系统的能级。
\end{enumerate}


\banswer{
	\begin{enumerate}
		\item 令 $\vec n=(\sin\theta,0,\cos\theta)$，并把题给哈密顿量写为
		\[
			H=-\Gamma\,\vec n\cdot\vec\sigma,\qquad
			\Gamma=\mu_B\hbar B_0
		\]
		（若题中 $S$ 已取为无量纲 $\sigma/2$，则改取
		$\Gamma=\mu_BB_0$）。于是
		\[
			U(t)=\cos\frac{\Gamma t}{\hbar}\,I
			+i\sin\frac{\Gamma t}{\hbar}\,\vec n\cdot\vec\sigma.
		\]
		记 $c=\cos(\Gamma t/\hbar)$、$s=\sin(\Gamma t/\hbar)$，作用于题给归一化初态得
		\[
			\boxed{\chi(t)=\frac13\begin{pmatrix}
			c(1-2i)+is[\cos\theta(1-2i)+2\sin\theta]\\[3pt]
			2c+is[\sin\theta(1-2i)-2\cos\theta]
			\end{pmatrix}}.
		\]

		\item 在基底
		$|\uparrow\uparrow\rangle,|\uparrow\downarrow\rangle,
		|\downarrow\uparrow\rangle,|\downarrow\downarrow\rangle$
		中，平行自旋态已经是本征态：
		\[
			\boxed{E_{\uparrow\uparrow}=a+b+c_0,\qquad
			E_{\downarrow\downarrow}=-a-b+c_0}.
		\]
		$m_z=0$ 子空间的矩阵为
		\[
			\begin{pmatrix}
			a-b-c_0&2c_0\\
			2c_0&-a+b-c_0
			\end{pmatrix},
		\]
		故另外两个能级为
		\[
			\boxed{E_\pm=-c_0\pm\sqrt{(a-b)^2+4c_0^2}}.
		\]
	\end{enumerate}
}



\item 
（共 30 分）考虑谐振子问题。
\begin{enumerate}
	%\renewcommand{\labelenumi}{\arabic{enumi}.}
	% A(\Alph) a(\alph) I(\Roman) i(\roman) 1(\arabic)
	%设定全局标号series=example	%引用全局变量resume=example
	%[topsep=-0.3em,parsep=-0.3em,itemsep=-0.3em,partopsep=-0.3em]
	%可使用leftmargin调整列表环境左边的空白长度 [leftmargin=0em]
	\item
	一维谐振子的哈密顿量为 $\hat{H}=\frac{\hat{p}^{2}}{2 m}+\frac{1}{2} k \hat{x}^{2}$，
	证明由不确定性关系得到的能量最
	小值与该谐振子的基态能量一致。
	\item 
	若(1)中的基态波函数是高斯型 $e^{-\beta x^{2}}$, 用变分法求 $\beta$ 。
	\item 
	利用升、降算符写出(1)中的第一激发态的波函数（不必归一)。
	\item 
	对于 三维 各 向 同 性 谐 振子 ，第一 激发 态是三 重简并 的。现 有一微 扰
	\[ 
	\hat{H}^{\prime}=b \hat{x} \hat{y}=\frac{\hbar b}{2 m \omega}\left(\begin{array}{ccc}0 & 1 & 0 \\ 1 & 0 & 0 \\ 0 & 0 & 0\end{array}\right)
	 \]
	写出该微扰引起的第一激发态的能级分裂。
\end{enumerate}


\banswer{
	令 $\omega=\sqrt{k/m}$。
	\begin{enumerate}
		\item 记 $d=\Delta x$，基态可取 $\langle x\rangle=\langle p\rangle=0$。由
		$\Delta p\ge\hbar/(2d)$，
		\[
			E\ge\frac{\hbar^2}{8md^2}+\frac12m\omega^2d^2.
		\]
		对 $d$ 极小化得到 $d^2=\hbar/(2m\omega)$ 及
		\[
			\boxed{E_{\min}=\frac12\hbar\omega},
		\]
		恰为谐振子精确基态能量。

		\item 对 $\psi=N e^{-\beta x^2}$，
		\[
			\langle x^2\rangle=\frac1{4\beta},\qquad
			\langle p^2\rangle=\hbar^2\beta.
		\]
		故
		\[
			E(\beta)=\frac{\hbar^2\beta}{2m}+\frac{k}{8\beta}.
		\]
		由 $dE/d\beta=0$ 得
		\[
			\boxed{\beta=\frac{m\omega}{2\hbar}
			=\frac{\sqrt{mk}}{2\hbar}}.
		\]

		\item
		\[
			\psi_1(x)\propto a^\dagger\psi_0(x)
			\propto\boxed{x\exp\left(-\frac{m\omega x^2}{2\hbar}\right)}.
		\]

		\item 在无扰第一激发态基
		$|100\rangle,|010\rangle,|001\rangle$ 中，题给微扰矩阵的本征值为
		\[
			\frac{\hbar b}{2m\omega},\qquad
			-\frac{\hbar b}{2m\omega},\qquad0.
		\]
		因此原来的 $5\hbar\omega/2$ 三重能级分裂为
		\[
			\boxed{\frac52\hbar\omega+\frac{\hbar b}{2m\omega},\quad
			\frac52\hbar\omega-\frac{\hbar b}{2m\omega},\quad
			\frac52\hbar\omega}.
		\]
		相应零级态为 $(|100\rangle\pm|010\rangle)/\sqrt2$ 与 $|001\rangle$。
	\end{enumerate}
}


\end{enumerate}

